\documentclass[10pt,landscape,a4paper]{article}
\usepackage[landscape,margin=1.2cm]{geometry}
\usepackage{multicol}
\usepackage{tabularx}
\usepackage{xcolor}
\usepackage{hyperref}
\usepackage[T1]{fontenc}
\usepackage{lmodern}
\usepackage{amssymb}
\usepackage{fuzz}
\usepackage{proof}

% Compact spacing
\setlength{\parindent}{0pt}
\setlength{\parskip}{2pt plus 1pt}
\setlength{\columnsep}{1.2cm}
\pagestyle{empty}

% Section formatting
\usepackage{titlesec}
\titleformat{\section}{\large\bfseries\color{darkgray}}{}{0pt}{}[\vspace{-4pt}\rule{\columnwidth}{0.4pt}]
\titlespacing{\section}{0pt}{6pt}{2pt}

% Code style
\newcommand{\kw}[1]{\texttt{\bfseries #1}}
\newcommand{\op}[1]{\texttt{#1}}
\newcommand{\renders}{\,$\rightarrow$\,}

% Stacked example: input then rendered output (compact)
\newenvironment{exwrite}%
  {\par\noindent\scriptsize\textbf{You write:}\par\vspace{0pt}\scriptsize}%
  {\vspace{1pt}}
\newenvironment{exget}%
  {\par\noindent\scriptsize\textbf{You get:}\par\vspace{0pt}\small}%
  {\par\vspace{3pt}}

% Table helper — left col monospace, right col roman
\newcolumntype{L}{>{\ttfamily}l}
\newcolumntype{R}{>{\raggedright\arraybackslash}X}

\begin{document}

\begin{center}
{\LARGE\bfseries txt2tex Relational Databases Reference}\\[2pt]
{\small Relational databases: primary keys, relational algebra, bindings, GROUP/UNGROUP.}\\[1pt]
{\footnotesize\color{gray}\url{https://github.com/jmf-pobox/txt2tex}}
\end{center}

\vspace{-2pt}

\begin{multicols}{3}
\raggedcolumns

% cheatsheet-common.tex — shared fragment for course-specific cheatsheet drivers.
% No preamble, no \begin{document}, no \begin{multicols}.
% Drivers must define: \kw, \op, \renders, exwrite/exget environments, L and R column types.
% This file is \input{} inside a \begin{multicols}{3} block.

% ── DOCUMENT STRUCTURE ──────────────────────────
\section{Document Structure}

\begin{tabularx}{\columnwidth}{LR}
=== Title ===        & \verb|\section*{Title}| \\
** Solution N **     & Bold heading with spacing \\
(a) Text             & Part label \\
TEXT: paragraph       & Prose block (smart formula detection) \\
LATEX: \textbackslash cmd & Raw \LaTeX\ passthrough \\
PAGEBREAK:           & Page break \\
LINEBREAK:           & Vertical gap (\verb|\medskip|) — between blocks, no page break \\
\end{tabularx}

\smallskip
\begin{tabularx}{\columnwidth}{LR}
TITLE: ...           & Document title \\
AUTHOR: ...          & Author name \\
DATE: ...            & Date string \\
BIBLIOGRAPHY: f.bib  & BibTeX file \\
{[cite key]}         & \verb|\citep{key}| \\
\end{tabularx}

% ── LOGIC ───────────────────────────────────────
\section{Logic}

\begin{tabularx}{\columnwidth}{LR}
p land q             & $p \land q$ \\
p lor q              & $p \lor q$ \\
lnot p               & $\lnot p$ \\
p => q               & $p \Rightarrow q$ \\
p <=> q              & $p \Leftrightarrow q$ \\
forall x : S | P     & $\forall x : S @ P$ \\
exists x : S | P     & $\exists x : S @ P$ \\
exists\_1 x : S | P  & $\exists_1 x : S @ P$ \\
lambda x : S | E     & $\lambda x : S @ E$ \\
mu x : S | P         & $\mu x : S @ P$ \\
\end{tabularx}

\smallskip
{\footnotesize\textbf{Note:} Use \kw{land}, \kw{lor}, \kw{lnot}. English \texttt{and}/\texttt{or}/\texttt{not} are not supported.}

% ── SETS ────────────────────────────────────────
\section{Sets \& Types}

\begin{tabularx}{\columnwidth}{LR}
x elem S             & $x \in S$ \\
x notin S            & $x \notin S$ \\
A subset B           & $A \subseteq B$ \\
A union B            & $A \cup B$ \\
A inter B            & $A \cap B$ \\
A setminus B         & $A \setminus B$ \\
power S              & $\power S$ \\
F S                  & $\finset S$ \\
\{x : S | P\}       & Set comprehension \\
A cross B            & $A \cross B$ \\
N / Z / N1           & $\nat$ / $\num$ / $\nat_1$ \\
n..m                 & $n \upto m$ \\
\end{tabularx}

% ── Z NOTATION BLOCKS ───────────────────────────
\section{Z Notation Blocks}

\begin{tabularx}{\columnwidth}{LR}
given A, B           & Given types: $[A, B]$ \\
T ::= a | b          & Free type \\
name == expr         & Abbreviation \\
\end{tabularx}

\smallskip
\begin{tabularx}{\columnwidth}{LR}
\multicolumn{2}{l}{\texttt{axdef} / \texttt{schema Name} / \texttt{gendef [X]}}\\
\quad declarations   & Variable declarations \\
\texttt{where}       & Separator \\
\quad predicates     & Constraining predicates \\
\texttt{end}         & Close the block \\
\end{tabularx}

% ── USAGE ───────────────────────────────────────
\section{Usage}

\begin{tabularx}{\columnwidth}{LR}
txt2tex f.txt        & $\rightarrow$ f.tex + f.pdf \\
txt2tex f.txt {-}{-}tex-only & $\rightarrow$ f.tex only \\
txt2tex -i           & Interactive REPL \\
txt2tex {-}{-}check-env  & Verify dependencies \\
\end{tabularx}


\end{multicols}

\newpage

\begin{center}
{\LARGE\bfseries txt2tex Relational Databases --- By Example}\\[2pt]
{\small Primary keys, relational algebra, bindings, GROUP/UNGROUP.}
\end{center}

\vspace{-2pt}

\begin{multicols}{3}
\raggedcolumns

% ── PRIMARY KEYS ────────────────────────────────
\section{Primary Keys}

{\scriptsize Mark a primary-key attribute with \kw{pk} in a named schema body. Fuzz-compatible.}

\begin{exwrite}
\begin{verbatim}
schema Class
  pk class : ClassName
  country  : CountryName
  bore     : N
end
\end{verbatim}
\end{exwrite}
\begin{exget}
\begin{center}
\begin{schema}{Class}
class : ClassName \\
country : CountryName \\
bore : \nat
\end{schema}\\[2pt]
\noindent$\mathrm{PK}(\mathrm{Class}) = \{class\}$
\end{center}
\end{exget}

{\scriptsize Composite key: two \kw{pk} lines $\Rightarrow$ $\mathrm{PK}(S)=\{a,b\}$.
Comma-separated names on one \kw{pk} line also work.
\kw{pk} is only valid inside a named \texttt{schema}; not allowed in \texttt{axdef}, \texttt{gendef}, or inline schema text.}

% ── RELATIONAL ALGEBRA ──────────────────────────
\section{Relational Algebra}

\begin{tabularx}{\columnwidth}{LR}
sigma[p](R)          & $\mathrm{Restrict}_{p}(R)$ — restrict \\
pi[A,B](R)           & $\mathrm{Project}\{A,B\}(R)$ — project \\
rho[A as B](R)       & $\mathrm{Rename}_{A \to B}(R)$ — rename \\
R bowtie S           & $R \otimes S$ — natural join \\
R bowtie [p] S       & $\mathrm{Join}_{p}(R, S)$ — theta-join \\
R div S              & $R \div S$ — division \\
\end{tabularx}

\smallskip
{\scriptsize\textbf{Naming a query.} Use \texttt{Name == expr} for any right-hand side — pure Z or relational algebra.}

\smallskip
{\scriptsize\textbf{Fuzz-compatible:} write algebra as free-standing lines in the document, not inside \texttt{zed}/\texttt{axdef}/\texttt{schema} blocks. Schemas and axdefs in the same document still type-check.}

\smallskip
{\scriptsize\textbf{Long expressions:} break at any operator — bare newline or trailing \texttt{\textbackslash}. Break \emph{before} \kw{setminus}, not after.}

\begin{exwrite}
\begin{verbatim}
BigGuns == pi[class,country](
  sigma[bore >= 16](Class))
\end{verbatim}
\end{exwrite}
\begin{exget}
\begin{center}
\noindent$BigGuns \defs \mathrm{Project}\{class,country\}(\mathrm{Restrict}_{bore \geq 16}(Class))$
\end{center}
\end{exget}

% ── BINDINGS ────────────────────────────────────
\section{Z Bindings \& Set Comprehension}

{\scriptsize Z RM §3.7 binding brackets construct labelled tuples.}

\begin{tabularx}{\columnwidth}{LR}
\{| name == e |\}    & $\lblot name == e \rblot$ \\
\{| a == e, b == f |\} & multiple components \\
\{| |\}              & empty binding \\
\end{tabularx}

\smallskip
{\scriptsize Multi-typed comprehension with binding:}

\begin{exwrite}
\begin{verbatim}
{ s : Ship; c : Class |
  s.class = c.class .
  {| name == s.name |} }
\end{verbatim}
\end{exwrite}
\begin{exget}
\begin{center}
$\{ s : Ship; c : Class \mid s.class = c.class \bullet \lblot name == s.name \rblot \}$
\end{center}
\end{exget}

{\scriptsize Use \texttt{;} to separate variable-type pairs inside \texttt{\{ \}}. Binding components are \textbf{comma}-separated (Z RM §3.7).}

\smallskip
{\scriptsize\textbf{Fuzz-compatible:} write bindings as free-standing expressions in the document, not inside \texttt{zed}/\texttt{axdef}/\texttt{schema} blocks.}

% ── GROUP / UNGROUP ──────────────────────────────
\section{GROUP \& UNGROUP}

{\scriptsize Date's nested-relation operators.}

\begin{tabularx}{\columnwidth}{LR}
R group (\{A,B\} as alias) & $R\ \mathop{\mathrm{GROUP}} (\{A,B\}\ \mathop{\mathrm{AS}}\ alias)$ \\
R ungroup alias    & $R\ \mathop{\mathrm{UNGROUP}}\ alias$ \\
\end{tabularx}

\begin{exwrite}
\begin{verbatim}
Members group ({name,email} as contact)
Members ungroup contact
\end{verbatim}
\end{exwrite}
\begin{exget}
\begin{center}
$Members\ \mathop{\mathrm{GROUP}}\ (\{name,email\}\ \mathop{\mathrm{AS}}\ contact)$\\[3pt]
$Members\ \mathop{\mathrm{UNGROUP}}\ contact$
\end{center}
\end{exget}

{\scriptsize\textbf{Fuzz-compatible:} write GROUP/UNGROUP as free-standing expressions in the document, not inside \texttt{zed}/\texttt{axdef}/\texttt{schema} blocks.}

% ── QUICK REFERENCE ─────────────────────────────
\section{Quick Reference}

\begin{tabularx}{\columnwidth}{LR}
pk a : T             & primary key in schema \\
pk a, b : T          & composite pk \\
sigma[p](R)          & restrict \\
pi[A](R)             & project \\
rho[A as B](R)       & rename attr \\
R bowtie S           & natural join ($\otimes$) \\
R div S              & division \\
\{| a == e |\}       & binding bracket \\
\{S; C | P . E\}    & multi-typed set comp \\
R group (\{A\} as x) & group attrs \\
R ungroup x          & ungroup attr \\
\end{tabularx}

% ── IDENTIFIERS ─────────────────────────────────
\section{Identifier Decoration}

\begin{tabularx}{\columnwidth}{LR}
count'               & $count'$ (after-state) \\
in?                  & $in?$ (input) \\
out!                 & $out!$ (output) \\
x?'                  & $x?'$ (runs may combine) \\
'sunk'               & $\mbox{`sunk'}$ (string literal) \\
\end{tabularx}

{\footnotesize Decoration attaches to any identifier. Reserved words (\texttt{theta}, \texttt{defs}, \texttt{hide}, \texttt{project}, \texttt{Delta}, \texttt{Xi}) cannot be decorated.}

\end{multicols}

\end{document}
